\chapter*{چکیده}
\addcontentsline{toc}{chapter}{چکیده}
\markboth{چکیده}{چکیده}

\textbf{سلام} یک زبانِ برنامه‌نویسیِ نوین، ساده و دوزبانه است که به شما اجازه
می‌دهد برنامه‌های خود را \emph{هم به فارسی و هم به انگلیسی} بنویسید. در سلام
می‌توانید واژگانِ کلیدی مانند \code{تابع}، \code{اگر}، \code{برای} و \code{بازگشت}
را به فارسی به‌کار ببرید، یا اگر بخواهید همان برنامه را با \code{func}، \code{if}،
\code{for} و \code{ret} به انگلیسی بنویسید. هر دو شکل، به یک برنامهٔ یکسان تبدیل
می‌شوند و با همان سرعت اجرا می‌شوند.

\section*{چرا سلام دوزبانه است؟}

سال‌هاست که آموزشِ برنامه‌نویسی برای فارسی‌زبانان با یک مانعِ پنهان همراه بوده
است: زبانِ انگلیسی. کودکی که هنوز انگلیسی نیاموخته، یا دانش‌آموزی که می‌خواهد
نخستین گام‌ها را بردارد، پیش از آنکه با \emph{منطقِ برنامه‌نویسی} آشنا شود،
ناچار است با واژگانِ بیگانه دست‌وپنجه نرم کند. سلام این مانع را برمی‌دارد. شما
می‌توانید با زبانِ مادری خود بیندیشید و بنویسید:

\begin{salam}
تابع اصلی:
    چاپ "سلام، دنیا!"
تمام
\end{salam}

در همان حال، سلام پلی به جهانِ گستردهٔ برنامه‌نویسیِ بین‌المللی نیز هست. چون هر
واژهٔ فارسی یک معادلِ دقیقِ انگلیسی دارد، خواننده به‌تدریج و بدونِ فشار، واژگانِ
استانداردِ جهانی را نیز می‌آموزد و آماده می‌شود تا در آینده با هر زبانِ دیگری کار
کند.

\section*{این کتاب برای چه کسانی است؟}

این کتاب کوشیده است \emph{همه} را در بر بگیرد:

\begin{itemize}
  \item \textbf{کودکانِ زیرِ ده سال:} با مثال‌های بازی‌گونه مانند کشیدنِ
        مثلثِ ستاره، جدولِ ضرب و بازیِ فیزباز، گام‌به‌گام و با زبانی ساده.
  \item \textbf{دانش‌آموزان:} برای یادگیریِ مفاهیمِ پایهٔ برنامه‌نویسی،
        الگوریتم‌ها و حلِ مسئله.
  \item \textbf{دانشجویان:} برای درکِ ساختارهای داده، شیءگرایی، کلی‌سازی
        و طراحیِ نرم‌افزار.
  \item \textbf{برنامه‌نویسانِ حرفه‌ای:} برای کار با ویژگی‌های پیشرفته مانند
        واسط‌ها، بستارها، هم‌روندی، ارتباط با کدِ خارجی و توسعهٔ وب.
\end{itemize}

\section*{چگونه این کتاب را بخوانیم؟}

اگر تازه‌کار هستید، کتاب را از فصلِ نخست و به ترتیب بخوانید؛ هر فصل بر فصلِ پیش
استوار است. اگر تجربه‌ای پیشین دارید، می‌توانید مستقیم به فصلِ موردِ نظر بروید.
در سراسرِ کتاب با این نشانه‌ها روبه‌رو می‌شوید:

\begin{note}
نکته‌ای برای روشن‌تر شدنِ مطلب یا یادآوریِ چیزی مهم.
\end{note}

\begin{tip}
ترفند یا روشی که کارِ شما را آسان‌تر یا برنامه‌تان را بهتر می‌کند.
\end{tip}

\begin{warn}
هشدار دربارهٔ اشتباهی رایج که بهتر است از آن دوری کنید.
\end{warn}

\begin{deep}
بخشی ژرف‌تر دربارهٔ آنکه \emph{چرا} و \emph{چگونه} چیزی در درونِ زبان کار
می‌کند. خواندنِ این بخش‌ها اختیاری است؛ اگر تازه‌کارید می‌توانید از آن‌ها بگذرید
و بعدها بازگردید.
\end{deep}

در پایانِ هر فصل، چند \textbf{تمرین} آمده است. پاسخِ همهٔ تمرین‌ها در پیوستِ
انتهای کتاب گردآوری شده است؛ اما پیش از دیدنِ پاسخ، خودتان بکوشید. برنامه‌نویسی
با \emph{انجام دادن} آموخته می‌شود، نه فقط با خواندن.

\vspace{0.5cm}
\begin{center}
\emph{به دنیای سلام خوش آمدید. بیایید نخستین برنامه را با هم بنویسیم.}
\end{center}
\clearpage
